\scanpage{645}
\mspage{33}

3. \uline{Cas des catégories à cofibrations}, \uline{des cat. à fibrations}.

Je rappelle les axiomes d'une \uline{structure à cofibrations}
$(W,C)$ sur une catégorie $M$.

$\mathrm{Cof}_0$ : $W$ est un localiseur.

$\mathrm{Cof}_1$ : $C$ contient les iso, est stable par \unsure{cochangemt}
de base, par composition.

$\mathrm{Cof}_2$ : $C\cap W$ est stable par \unsure{cochangement} de base
(et aussi par composition, et contient les
iso, en vertu de $\mathrm{Cof}_0,\mathrm{Cof}_1$)

$\mathrm{Cof}_3$ : Toute flèche $f$ \unsure{dans|de} $M$ se factorise en
$f=pi$, avec $i\in C,p\in W$.

(Axiomes de K. S. Brown, légèrement généralisés
par Anderson.) Il y a un ``axiome de
continuité'' \inserted{ou ``\uline{\unsure{axiome des fins|axiome des feux|axiome sans fin}}''} facultatif, qui est le suivant

$\mathrm{Cof}_4$ : \inserted{$M$ \unsure{stable} par petites $\varinjlim$ filtrantes, et}
$C$ et $TC\overset{\mathrm{def}}{=}C\cap W$ sont stables par
composition transfinie. \deleted{\unsure{ou}}.

Ainsi, $\mathrm{Cof}_1$, $\mathrm{Cof}_2$, $\mathrm{Cof}_4$ peuvent se résumer
en disant que $C$ et $TC$ sont $Q$-saturés
\unsure{transfiniment} à gauche.

La structure de \unsure{cofibration|cofibrations} est dite
\uline{spéciale} si elle, \deleted{est stable} satisfait $\mathrm{Cof}_4$,
si $M$ est spéciale, et si $C$ admet
une \uline{petite} partie $C_0$ qui soit $Q$-génératrice
à gauche, i.e. telle que

\[
 C\subset \unsure{C_0^{\sim}|\widetilde{C_0}}\quad\text{i.e.}\quad\widetilde C=C_0^{\sim}.\tag*{(3.1)}
\]
% archive sheet number: 644
